\chapter{Einleitung}
\label{chap:expose}

\section{Kontext und Motivation}

Die Digitalisierung von Geschäftsprozessen und die steigende Verlagerung von IT-Infrastrukturen in die Cloud haben das \ac{IAM} zu einem zentralen Sicherheitsbaustein der Unternehmensarchitektur gemacht. Organisationen nutzen im Durchschnitt 3,4 verschiedene Cloud-Anbieter\footnote{\label{fn:sqmagazine}\url{https://sqmagazine.co.uk/cloud-adoption-statistics}, aufgerufen am 13.10.2025}, wobei 753 befragte Unternehmen angeben, dass 89\% der Unternehmen eine Multi-Cloud-Strategie verfolgen\footnote{\url{https://www.flexera.com/blog/finops/cloud-computing-trends-flexera-2024-state-of-the-cloud-report/}, aufgerufen am 13.10.2025}. Beispielsweise werden die Vorteile der verschiedenen Anbieter kombiniert, um Kosten und Performance zu optimieren. Diese Multi-Cloud-Umgebungen bringen neue Herausforderungen für die IT-Sicherheit von Identitäts- und Zugriffsmanagementsystemen mit sich. Zudem ist die IT-Sicherheit an zweiter Stelle bei den größten Herausforderungen von Unternehmen in der Cloud\footref{fn:sqmagazine}.
\newline
Beim Entwickeln von Web-Anwendungen mit Backend-Systemen in der Cloud müssen Entwickler entscheiden, wie sie Benutzer authentifizieren und deren Berechtigungen verwalten. Cloud-Anbieter bieten dafür verschiedene Lösungen an, beispielsweise bietet \ac{AWS} Cognito\footnote{\label{fn:whatcognito}\url{https://docs.aws.amazon.com/cognito/latest/developerguide/what-is-amazon-cognito.html}, aufgerufen am 13.10.2025} an und SAP setzt auf \ac{SAP IAS}\footnote{\label{fn:sapias}\url{https://www.sap.com/germany/products/technology-platform/identity-authentication.html}, aufgerufen am 16.10.2025} in Kombination mit dem \ac{SAP CAP}\footnote{\label{fn:sapcap}\url{https://cap.cloud.sap/docs/guides/security/}, aufgerufen am 13.10.2025}. Diese Systeme unterstützen etablierte Standards wie \textit{OAuth 2.0} (Autorisierungsframework für API-Zugriff)\footnote{\url{https://datatracker.ietf.org/doc/html/rfc6749}, aufgerufen am 16.10.2025}, \textit{OpenID Connect} (Authentifizierungsschicht auf OAuth 2.0-Basis)\footnote{\url{https://openid.net/specs/openid-connect-core-1_0.html}, aufgerufen am 16.10.2025} oder \textit{SAML 2.0} (XML-basierter Standard für Single Sign-On)\footnote{\url{https://docs.oasis-open.org/security/saml/Post2.0/sstc-saml-tech-overview-2.0.html}, aufgerufen am 16.10.2025}, implementieren diese jedoch unterschiedlich \citep{Ali2025}.
\newline
AWS Cognito ist eine spezialisierte \ac{CIAM}-Lösung von \ac{AWS}. Sie besteht aus zwei Hauptkomponenten: User Pools für die Authentifizierung und Autorisierung von Anwendungsbenutzern, und Identity Pools für den sicheren Zugriff auf \ac{AWS}-Services\footnote{\ref{fn:whatcognito}}. AWS Cognito unterstützt OAuth 2.0, OpenID Connect und SAML 2.0 und bietet Features wie Multi-Faktor-Authentifizierung, Social Login und eine anpassbare Benutzeroberfläche\footnote{\label{fn:cognito}\url{https://aws.amazon.com/cognito/features/}, aufgerufen am 13.10.2025}.
\newline
\ac{SAP CAP} ist ein Framework für die Entwicklung von Unternehmensanwendungen auf der \ac{SAP BTP}. Für das Identitätsmanagement integriert \ac{SAP CAP} mit dem \ac{SAP IAS}, einem OpenID Connect-kompatiblen Service. \ac{SAP IAS} übernimmt die zentrale Authentifizierung der Nutzer, während die Autorisierung durch die Applikation verwaltet wird\footnote{\ref{fn:sapcap}}.
\newline
Laut der \ac{OWASP} Cloud-Native Application Security Top 10 ist „Improper authentication \& authorization"\ das dritthäufigste Sicherheitsrisiko bei Cloud-Anwendungen\footnote{\label{fn:owaspcloudnative}\url{https://owasp.org/www-project-cloud-native-application-security-top-10/}, aufgerufen am 13.10.2025}. Häufige Probleme sind dabei fehlende Authentifizierung bei APIs, zu weitreichende Berechtigungen oder unsichere Kommunikation zwischen verschiedenen Services. Deshalb ist es wichtig zu verstehen, wie die verschiedenen Cloud-Anbieter diese Risiken adressieren und welche Unterschiede in der praktischen Umsetzung bestehen.
\newline
In der Praxis zeigen sich zudem verschiedene Herausforderungen. Zu einem unterscheidet sich die Konfigurationskomplexität zwischen den Anbietern. Während AWS Cognito eine große Varianz von Konfigurationsmöglichkeiten bietet, kann dies unter anderem zu Fehlkonfigurationen führen, die z.B. Verwundbarkeiten verursachen\footnote{\url{https://docs.aws.amazon.com/cognito/latest/developerguide/cognito-user-pool-security.html}, aufgerufen am 16.10.2025}. \ac{SAP IAS} setzt dabei auf \textit{Anwendungsintegration}, die unabhängige Softwareanwendungen zu einer funktionierenden Einheit verbindet\footnote{\url{https://www.sap.com/germany/products/technology-platform/what-is-enterprise-integration/application-integration.html}, aufgerufen am 16.10.2025}, was die Einrichtung für kleinere Anwendungen komplexer gestalten kann\footnote{\url{https://help.sap.com/docs/identity-authentication/identity-authentication/configuring-applications}, aufgerufen am 16.10.2025}. Zusätzlich variieren die Konfigurationen, denn AWS Cognito User Pools bieten beispielsweise verschiedene Passwort-Policy-Optionen und MFA-Konfigurationen\footnote{\url{https://docs.aws.amazon.com/cognito/latest/developerguide/user-pool-settings-policies.html}, aufgerufen am 16.10.2025}, während \ac{SAP IAS} als zentraler \ac{IdP} für SAP-Landschaften entwickelt wurde und daher standardmäßig auf Unternehmensintegration ausgerichtet ist und auch die Anbindung an bestehende SAP-Systeme, vordefinierte Rollen und Berechtigungskonzepte unterstützt\footnote{\ref{fn:sapias}}. Diese unterschiedlichen Ansätze können während der Entwicklung zu verschiedenen Implementierungsaufwänden, unterschiedlichen Ausprägungen von Entwickler-Usability (Setup-Zeit, API-Design, Dokumentationsqualität und Konfigurationskomplexität) und IT-Sicherheit (Authentifizierungsmechanismen, Autorisierungskonzepte, Fehlkonfigurationen) sowie variierender Wartbarkeit führen.

\section{Ziele}
Das Ziel dieser Arbeit ist die systematische Evaluation und der Vergleich von AWS Cognito und \ac{SAP CAP} mit \ac{SAP IAS} als Cloud-\ac{IdP} . Dazu wird eine \textit{evaluative comparative Case Study} nach \citep{Kitchenham1995} durchgeführt, die beide Systeme anhand definierter Hypothesen und Evaluationskriterien wie Sicherheitsfeatures, Entwicklerorientiertheit und Integration systematisch analysiert und vergleicht. Basierend auf den Ergebnissen werden praxisorientierte Empfehlungen und ein Entscheidungsleitfaden für die Auswahl geeigneter \ac{IdP} entwickelt. Zudem werden Best Practices für beide Identitäts- und Zugriffsmanagementsysteme erstellt, die die zuvor erworbenen Erkenntnisse umsetzen.
\begin{enumerate}
    \item[\namedlabel{z1}{Ziel 1}] \textbf{Theoretische Grundlagen schaffen}
    \begin{enumerate}
        \item[\namedlabel{z11}{Ziel 1.1}] OAuth 2.0, OpenID Connect und SAML sind erklärt und deren Umsetzung in AWS Cognito und \ac{SAP IAS} ist dokumentiert. Die Dokumentation ist fertig bis zum 10.11.2025.
        \item[\namedlabel{z12}{Ziel 1.2}] Allgemeine Cloud-Sicherheitsrisiken sind analysiert und deren Bedeutung für Identity-Management ist bewertet mit der Einbeziehung von \ac{OWASP}. Die Analyse ist abgeschlossen bis zum 14.11.2025.
    \end{enumerate}
    
    \item[\namedlabel{z2}{Ziel 2}] \textbf{Einführung in Methodik mit Hypothesen und Identity-Provider analysieren}
    \begin{enumerate}
        \item[\namedlabel{z21}{Ziel 2.1}] Messbare Bewertungskriterien sind definiert (z.B. Wartbarkeit, Programming Understanding, Reliability, Benutzerfreundlichkeit) bis zum 17.11.2025.
        \item[\namedlabel{z22}{Ziel 2.2}] AWS Cognito Funktionen und Architektur sind detailliert untersucht und dokumentiert. Die Analyse ist fertig bis zum 20.11.2025.
        \item[\namedlabel{z23}{Ziel 2.3}] Einführung in die empirische Methodik und Hypothesen sind basierend auf den vorherigen Erkenntnissen bis zum 25.11.2025 aufgestellt.
        \item[\namedlabel{z24}{Ziel 2.4}] \ac{SAP IAS} mit \ac{SAP CAP} ist in gleicher Weise analysiert und dokumentiert. Die Untersuchung ist abgeschlossen bis zum 30.11.2025.
        \item[\namedlabel{z25}{Ziel 2.5}] Beide Systeme sind nach den gleichen Kriterien bewertet und verglichen. Der Vergleich ist fertig bis zum 05.12.2025.
    \end{enumerate}
    
    \item[\namedlabel{z3}{Ziel 3}] \textbf{Praktische Implementierung durchführen}
    \begin{enumerate}
        \item[\namedlabel{z31}{Ziel 3.1}] Eine Web-Anwendung mit AWS Cognito ist implementiert (User Registrierung, Login, gesicherte Zugriffe über Pfade). Der Prototyp ist fertig bis zum 15.12.2025.
        \item[\namedlabel{z32}{Ziel 3.2}] Die gleiche Anwendung mit \ac{SAP CAP} und \ac{SAP IAS} ist implementiert. Der Prototyp ist fertig bis zum 25.12.2025.
        \item[\namedlabel{z33}{Ziel 3.3}] Die Gebrauchstauglichkeit und
        praktische Nutzbarkeit beider Varianten aus Entwicklersicht sind
        systematisch evaluiert und dokumentiert. Die Evaluation umfasst
        ausgearbeitete Kriterien wie Implementierungsaufwand,
        Dokumentationsqualität, API-Usability und Anwendungsfreundlichkeit
        bis zum 30.12.2025.
    \end{enumerate}
    
    \item[\namedlabel{z4}{Ziel 4}] \textbf{Hypothesen testen}
    \begin{enumerate}
        \item[\namedlabel{z41}{Ziel 4.1}] Sicherheitseigenschaften beider Systeme sind getestet und gemessen. Die Tests sind durchgeführt bis zum 05.01.2026.
        \item[\namedlabel{z42}{Ziel 4.2}] Bestimmte vorherige definierte Kriterien beider Systeme sind gemessen und verglichen bis zum 10.01.2026.
        \item[\namedlabel{z43}{Ziel 4.3}] Die aufgestellten Hypothesen sind durch die Messdaten bestätigt oder widerlegt. Die Auswertung ist abgeschlossen bis zum 15.01.2026.
    \end{enumerate}
    
    \item[\namedlabel{z5}{Ziel 5}] \textbf{Ergebnisse auswerten und Empfehlungen ableiten}
    \begin{enumerate}
        \item[\namedlabel{z51}{Ziel 5.1}] Ein Leitfaden für die Auswahl zwischen beiden Identity-Providern ist erstellt. Der Leitfaden ist fertig bis zum 20.01.2026.
        \item[\namedlabel{z52}{Ziel 5.2}] Best Practices für die sichere Implementierung sind dokumentiert. Die Dokumentation ist fertig bis zum 25.01.2026.
        \item[\namedlabel{z53}{Ziel 5.3}] Die Ergebnisse sind zusätzlich durch Literatur oder vorhanden Best-Practices validiert. Die Validierung ist abgeschlossen bis zum 31.01.2026.
    \end{enumerate}
\end{enumerate}

\section{Vorgehen}

\begin{description}
    \item \textbf{Vorgehen 1} Case Study Kontext definieren nach \citep{Kitchenham1995} (entspricht \ref{z11}, \ref{z12}, \ref{z21})
        \item \textbf{Vorgehen 1.1} Grundlagen werden durch Literatur erklärt und die Umsetzung in AWS Cognito und \ac{SAP CAP} ist angedeutet.
        \item \textbf{Vorgehen 1.2} Bewertungskriterien werden systematisch gesucht und mithilfe von \ac{OWASP} und anderer relevanter Literatur definiert.
        \item \textbf{Vorgehen 1.3} Die Ziele der Case Study werden klar definiert: Vergleich von AWS Cognito und \ac{SAP IAS} hinsichtlich erarbeiteter Kriterien, z.B. Sicherheit, Entwickler-Usability, Wartbarkeit.
    
    \item \textbf{Vorgehen 2} Hypothesen aufstellen und Identity-Provider analysieren (entspricht \ref{z22}, \ref{z23}, \ref{z24}, \ref{z25})
        \item \textbf{Vorgehen 2.1} Die Evaluationshypothesen werden formuliert (z.B. \glqq AWS Cognito ist einfacher zu implementieren als SAP IAS\grqq\
         oder \glqq SAP IAS bietet bessere Enterprise-Integration\grqq).
        \item \textbf{Vorgehen 2.2} Messbare Variablen werden definiert, wie z.B. Setup-Zeit, Lines of Code, Sicherheitsfeatures, Performance-Metriken.
        \item \textbf{Vorgehen 2.3} Beide Identity-Provider werden systematisch analysiert und charakterisiert und externe Projektbeschränkungen werden dokumentiert: Zeitrahmen, verfügbare Ressourcen und technische Einschränkungen.
    
    \item \textbf{Vorgehen 3} Case Study planen (entspricht \ref{z31}, \ref{z32}, \ref{z33})
        \item \textbf{Vorgehen 3.1} Experimentelle Subjekte und Objekte werden festgelegt: Standard-Web-Anwendung mit User Registration, Login und gesicherte Zugriffe über Pfade als Vergleichsbasis.
        \item \textbf{Vorgehen 3.2} Der Zeitpunkt der Methodennutzung wird definiert: Implementierung während der Entwicklungsphase.
        \item \textbf{Vorgehen 3.3} Der Zeitpunkt der Messungen wird festgelegt: Kontinuierliche Messung während der Implementierung und Tests nach Fertigstellung.
        \item \textbf{Vorgehen 3.4} Beide Prototypen werden unter kontrollierten Bedingungen implementiert und alle Metriken werden systematisch erfasst.
    
    \item \textbf{Vorgehen 4} Hypothesen validieren (entspricht \ref{z41}, \ref{z42}, \ref{z43})
        \item \textbf{Vorgehen 4.1} Die benötigten Daten werden gesammelt und die Messbarkeit wird sichergestellt.
        \item \textbf{Vorgehen 4.2} Mögliche „Effekte“ der beiden \ac{IdP} werden identifiziert und von anderen Einflüssen getrennt.
        \item \textbf{Vorgehen 4.3} Die korrekte Nutzung beider Methoden wird durch Dokumentation und Best Practices sichergestellt.
        \item \textbf{Vorgehen 4.4} Wichtige Zustandsvariablen werden identifiziert und kontrolliert (z.B. Entwicklererfahrung, Komplexität, ...).
        \item \textbf{Vorgehen 4.5} Performance-Tests, Security-Tests und Usability-Tests werden durchgeführt, um die Hypothesen zu prüfen.
    
    \item \textbf{Vorgehen 5} Ergebnisse analysieren (entspricht \ref{z51}, \ref{z52}, \ref{z53})
        \item \textbf{Vorgehen 5.1} Die Ergebnisse der Fallstudie werden nach statistischen Methoden analysiert und die Hypothesen werden bestätigt oder widerlegt.
        \item \textbf{Vorgehen 5.2} Die Evaluationsergebnisse werden bewertet und dokumentiert.
        \item \textbf{Vorgehen 5.3} Ein Entscheidungsleitfaden wird aus den validierten Ergebnissen entwickelt.
        \item \textbf{Vorgehen 5.4} Die Ergebnisse werden durch Literatur validiert und Best Practices werden abgeleitet.
\end{description}


\section{Wissenschaftlicher Beitrag}

Diese Arbeit leistet einen wissenschaftlichen Beitrag zur systematischen Bewertung von Cloud-Identity-Providern durch eine \textit{evaluative comparative Case Study} nach \citep{Kitchenham1995}. Dadurch wird eine hypothesenbasierte Evaluation der praktischen und theoretischen Implementierungsunterschiede zwischen führenden \ac{IdP}s dargestellt.
\newline
Der Beitrag umfasst drei wesentliche Aspekte: Zunächst wird ein systematisches Bewertungsframework entwickelt, das messbare Kriterien für z.B. Sicherheit, Entwicklerfreundlichkeit und Wartbarkeit definiert. Danach werden durch kontrollierte Prototypimplementierung Daten zum Implementierungsaufwand und zur praktischen Nutzbarkeit von AWS Cognito und \ac{SAP IAS} erhoben. Zuletzt werden die aufgestellten Hypothesen durch Tests validiert oder widerlegt.
\newline
Die Ergebnisse unterstützen Entwickler bei der Auswahl geeigneter \ac{IdP} und tragen zu der Forschung in der systematischen Evaluation von Cloud-\newline Authentifizierungslösungen bei. Der entwickelte Entscheidungsleitfaden und die validierten Best Practices bieten konkrete Hilfestellungen für die Praxis und dienen als Grundlage für weitere Forschungen oder Vergleiche in diesem Bereich.


